\documentclass[12pt, a4paper]{article}

% ─── Paquetes esenciales ──────────────────────────────────────────────────────
\usepackage[utf8]{inputenc}
\usepackage[T1]{fontenc}
\usepackage{geometry}
\geometry{margin=2.5cm, headheight=15pt}
\usepackage{setspace}
\onehalfspacing

\usepackage{graphicx}
\usepackage{hyperref}
\hypersetup{
    colorlinks=true,
    linkcolor=black,
    citecolor=blue,
    urlcolor=blue
}
\usepackage{amsmath, amssymb}
\usepackage{booktabs}
\usepackage{tabularx}
\usepackage{array}
\usepackage{longtable}
\usepackage{xcolor}
\usepackage{listings}
\usepackage{caption}
\usepackage{subcaption}
\usepackage{enumitem}
\usepackage{fancyhdr}
\usepackage{titlesec}
\usepackage{abstract}
\usepackage{cite}
% microtype removed for compatibility

% ─── Configuración de cabeceras ───────────────────────────────────────────────
\pagestyle{fancy}
\fancyhf{}
\rhead{\small Seguridad en Aplicaciones}
\lhead{\small Observabilidad e Inicio de Sesión Seguro}
\cfoot{\thepage}
\renewcommand{\headrulewidth}{0.4pt}

% ─── Estilo de secciones ──────────────────────────────────────────────────────
\titleformat{\section}{\large\bfseries}{\thesection.}{0.5em}{}
\titleformat{\subsection}{\normalsize\bfseries}{\thesubsection.}{0.5em}{}
\titleformat{\subsubsection}{\normalsize\itshape}{\thesubsubsection.}{0.5em}{}

% ─── Configuración de listings (código) ───────────────────────────────────────
\lstset{
  basicstyle=\small\ttfamily,
  breaklines=true,
  frame=single,
  backgroundcolor=\color{gray!8},
  keywordstyle=\color{blue}\bfseries,
  commentstyle=\color{gray},
  stringstyle=\color{red!70!black},
  numbers=left,
  numberstyle=\tiny\color{gray},
  xleftmargin=1.5em,
  framexleftmargin=1.5em,
  captionpos=b
}

% ═════════════════════════════════════════════════════════════════════════════
\begin{document}
% ═════════════════════════════════════════════════════════════════════════════

% ─── Portada ─────────────────────────────────────────────────────────────────
\begin{titlepage}
  \centering
  \vspace*{1.5cm}

  {\LARGE \textbf{De Contraseñas a Passkeys: Evolución del Inicio de Sesión
  Seguro y Observabilidad en Aplicaciones Web, Móviles y de Escritorio}}\\[1.8cm]

  {\large Trabajo de Grado --- Seguridad en Aplicaciones}\\[0.5cm]
  {\large Ingeniería de Sistemas}\\[2.5cm]

  {\normalsize \textit{Presentado por:}}\\[0.4cm]
  {\large \textbf{[Nombre del estudiante]}}\\[2cm]

  {\normalsize \textit{Docente:}}\\[0.4cm]
  {\large \textbf{[Nombre del docente]}}\\[2.5cm]

  {\large [Nombre de la Universidad]}\\[0.3cm]
  {\large Facultad de Ingeniería}\\[1.5cm]

  {\large \today}
  \vfill
\end{titlepage}

% ─── Resumen / Abstract ───────────────────────────────────────────────────────
\begin{abstract}
\noindent
La gestión de identidades digitales y la supervisión del comportamiento de las
aplicaciones representan dos de los desafíos más críticos de la ingeniería de
software contemporánea. Los esquemas de autenticación fundamentados en
contraseñas estáticas han demostrado ser sistémicamente vulnerables frente a
ataques de phishing, credential stuffing y toma de control de cuentas (ATO,
\textit{Account Takeover}), vectores responsables de pérdidas económicas
globales estimadas en 9,5 billones de dólares anuales \cite{authx2025}.
Paralelamente, la creciente complejidad de las arquitecturas distribuidas ---
microservicios, contenedores efímeros y despliegues multinube --- invalida los
modelos de monitoreo tradicional basados en umbrales estáticos, incapaces de
correlacionar eventos de alto nivel con acciones maliciosas a nivel de kernel.

Este trabajo presenta un análisis técnico y arquitectónico de los mecanismos
modernos de inicio de sesión seguro y de los paradigmas de observabilidad
profunda. Se examina la transición desde contraseñas hasta passkeys bajo los
estándares FIDO2/WebAuthn y el marco OAuth 2.0 con PKCE, diferenciando las
estrategias de implementación en plataformas web, móviles y de escritorio. Se
analiza además la integración de OpenTelemetry y eBPF (\textit{extended
Berkeley Packet Filter}) como pilares de observabilidad, contrastándolos con el
monitoreo clásico. El trabajo concluye con la evaluación de conformidad de la
arquitectura propuesta frente al estándar NIST SP 800-63-4 \cite{nist80063},
el OWASP Top 10 \cite{owasptop10} y el marco normativo colombiano
(Leyes 1273 de 2009 y 1581 de 2012).

\bigskip
\noindent\textbf{Palabras clave:} autenticación, passkeys, FIDO2, WebAuthn,
OAuth 2.0, PKCE, observabilidad, OpenTelemetry, eBPF, Zero Trust, seguridad
de aplicaciones, NIST SP 800-63-4.
\end{abstract}

\newpage
\tableofcontents
\newpage

% ═════════════════════════════════════════════════════════════════════════════
\section{Introducción}
% ═════════════════════════════════════════════════════════════════════════════

La transformación digital ha disuelto los perímetros de red tradicionales. La
proliferación del cómputo en la nube, los ecosistemas de dispositivos móviles, el
modelo BYOD (\textit{Bring Your Own Device}) y las arquitecturas de
microservicios han fragmentado la infraestructura hasta el punto en que cada
componente --- desde un contenedor efímero hasta un dispositivo personal ---
representa una superficie de ataque individual que debe protegerse de forma
independiente \cite{zerotrust2025}. En este entorno hiperconectado, la
\textbf{identidad del usuario} se ha convertido en la nueva frontera de control.

El inicio de sesión es el mecanismo algorítmico que determina quién interactúa
con el sistema, en qué contexto y con qué nivel de privilegios. Sin embargo, el
modelo históricamente dominante --- la contraseña estática como ``secreto
compartido'' --- presenta vulnerabilidades estructurales insalvables. La
reutilización masiva de credenciales (más del 50\% de los usuarios reutiliza la
misma contraseña en al menos cinco cuentas \cite{frontegg2025}), combinada con
el acceso industrializado a bases de datos filtradas en la Dark Web, ha
convertido el \textit{credential stuffing} y el ATO en fenómenos de escala
industrial.

De forma paralela, la incapacidad del monitoreo tradicional para detectar
intrusiones basadas en credenciales legítimas expone una brecha operacional
profunda. Cuando un atacante accede con credenciales válidas pero robadas, los
indicadores clásicos de salud del sistema (CPU, memoria, tráfico de red) no
registran anomalía alguna, lo cual es reconocido por OWASP como ``Fallas en el
Registro y Monitoreo de Seguridad'' (A09:2021) \cite{owasptop10}.

Este trabajo aborda la intersección de ambas problemáticas, proponiendo una
arquitectura de ``Seguridad por Diseño'' (\textit{Secure by Design})
\cite{owaspSBD} que integra (a) flujos de autenticación sin contraseña basados
en criptografía asimétrica y (b) una malla de observabilidad profunda
sustentada en OpenTelemetry y eBPF.

% ═════════════════════════════════════════════════════════════════════════════
\section{Marco Teórico}
% ═════════════════════════════════════════════════════════════════════════════

\subsection{Fundamentos de Seguridad de la Información}

La seguridad de la información descansa sobre la tríada CIA (\textit{Confidentiality,
Integrity, Availability}), a la que los marcos modernos agregan el
\textbf{No Repudio} como cuarto pilar \cite{iso27001}:

\begin{itemize}[noitemsep]
  \item \textbf{Confidencialidad:} Acceso restringido a entidades autorizadas,
        garantizado mediante cifrado de estado del arte (AES-256 en reposo,
        TLS 1.3 en tránsito).
  \item \textbf{Integridad:} Inmutabilidad de los datos, respaldada por
        funciones hash criptográficas (SHA-256 o superior) y firmas digitales.
  \item \textbf{Disponibilidad:} Acceso oportuno a recursos para usuarios
        legítimos, mitigando amenazas DoS/DDoS.
  \item \textbf{No Repudio:} Garantía criptográfica de que una acción no puede
        ser negada por quien la originó.
\end{itemize}

La estrategia de \textbf{Defensa en Profundidad} (\textit{Defense in Depth})
asume que ningún control individual es invulnerable, por lo que exige capas
superpuestas de protección \cite{iso27001}. La \textbf{Arquitectura Zero Trust}
\cite{zerotrust2025,nistZTA} materializa este principio con el axioma
``\textit{nunca confiar, siempre verificar}'': toda solicitud de acceso debe
autenticarse y autorizarse dinámicamente, independientemente de la ubicación
topológica del solicitante.

\subsection{Panorama de Amenazas}

El contexto global y regional justifica la urgencia de adoptar controles modernos.
A nivel mundial, los costos del cibercrimen superan los 9,5 billones de dólares
anuales \cite{authx2025}. En Colombia, el MinTIC reportó más de 36.000 millones
de intentos de ciberataque durante 2024 \cite{acis2024}, mientras las denuncias
por delitos informáticos crecieron un 23\% interanual, superando los 77.600 casos
\cite{ccit2024}. Los vectores más frecuentes continúan siendo el phishing dirigido,
el ransomware y la suplantación de identidad digital.

Los ataques de \textit{credential stuffing} ilustran la industrialización del
fraude: utilizando botnets e infraestructuras de proxies residenciales, los
actores maliciosos automatizan millones de intentos de autenticación con
credenciales filtradas. Análisis del tráfico corporativo revelan que en promedio
el 19\% de todos los intentos de autenticación diarios corresponde a ataques de
este tipo, cifra que puede alcanzar el 44\% durante picos coordinados
\cite{verizon2025}.

\subsection{Evolución de la Autenticación Digital}
\label{sec:evolucion}

La industria ha transitado por cuatro etapas evolutivas claramente diferenciadas:

\subsubsection{Contraseñas Estáticas (Secreto Compartido)}

El modelo tradicional almacena en el servidor un hash de la contraseña del usuario.
Su vulnerabilidad estructural radica en que ambos extremos del canal deben
conocer (o derivar) el mismo secreto. Una brecha en la base de datos del servidor
compromete todas las cuentas, especialmente cuando se usan algoritmos obsoletos
(MD5, SHA-1 sin \textit{salting}) \cite{googleModernPass}.

\subsubsection{Autenticación Multifactor Básica (MFA)}

La MFA clásica combina \textit{algo que el usuario sabe} (contraseña) con
\textit{algo que posee} (dispositivo móvil) o \textit{algo que es} (biometría).
Los códigos OTP por SMS, sin embargo, son susceptibles al secuestro del protocolo
SS7, la clonación de SIM (\textit{SIM swapping}) y los ataques
AiTM\footnote{AiTM: \textit{Adversary-in-the-Middle}, donde el atacante
intercepta en tiempo real la contraseña y el código OTP.} \cite{srivathsav2025}.

\subsubsection{MFA Resistente al Phishing}

Los estándares FIDO U2F introdujeron llaves de seguridad hardware cuya respuesta
criptográfica está vinculada al \textit{Origin} (nombre de dominio verificado) del
sitio. Si el usuario es engañado y navega a una URL falsa, el token se niega a
generar una firma válida para ese dominio, abortando el ataque
\cite{frontegg2025,srivathsav2025}.

\subsubsection{Passkeys y el Estándar FIDO2 / WebAuthn}
\label{sec:passkeys}

Los \textbf{passkeys} representan la eliminación total de la contraseña como factor
humano. Fundamentados en la especificación WebAuthn del W3C y el protocolo CTAP2
de la FIDO Alliance \cite{mojoauth2025}, operan mediante criptografía asimétrica:

\begin{enumerate}[noitemsep]
  \item \textbf{Registro:} El dispositivo genera un par de claves asimétricas
        único para cada dominio. La \textbf{clave privada} queda protegida en el
        enclave seguro del dispositivo (Secure Enclave en iOS/macOS, TPM en
        Windows); la \textbf{clave pública} se almacena en el servidor.
  \item \textbf{Autenticación:} El servidor emite un \textit{challenge} aleatorio.
        El dispositivo, previa verificación biométrica local, firma el challenge
        con la clave privada. El servidor verifica la firma con la clave pública
        registrada.
\end{enumerate}

\noindent La clave privada \textbf{jamás abandona el dispositivo} y está
bloqueada por biometría (FaceID, TouchID, Windows Hello) o PIN local. Este diseño
erradica algorítmicamente el \textit{credential stuffing}: no hay contraseñas
almacenadas en los servidores que puedan robarse. Empíricamente, los passkeys
logran tasas de éxito en el inicio de sesión superiores al 98\%, reducen la
fricción hasta en un 80\% y ofrecen inmunidad nativa al phishing
\cite{mojoauth2025,authsignal2025}.

\subsection{Arquitectura de Seguridad Moderna: OAuth 2.0, OIDC y JWT}
\label{sec:oauth}

En arquitecturas distribuidas, la autenticación se desacopla lógicamente de los
servicios de negocio mediante protocolos estándar:

\subsubsection{JSON Web Tokens (JWT)}

Tras validar las credenciales, el servidor de identidad emite un JWT firmado
criptográficamente. El token contiene \textit{claims} sobre la identidad y los
permisos del usuario en su \textit{payload}. Los microservicios validan
matemáticamente la autenticidad del token verificando su firma, sin consultas
sincrónicas a la base de datos central, lo que maximiza escalabilidad y
rendimiento \cite{frontegg2025}.

\subsubsection{OAuth 2.0 Authorization Code con PKCE}

Los clientes públicos (SPAs, aplicaciones móviles nativas, apps de escritorio)
no pueden almacenar de forma segura un \textit{client\_secret}. La extensión
\textbf{PKCE} (\textit{Proof Key for Code Exchange}) resuelve esta limitación
mediante inyección de entropía dinámica \cite{srivathsav2025}:

\begin{enumerate}[noitemsep]
  \item El cliente genera aleatoriamente un \texttt{code\_verifier} y calcula
        su hash SHA-256: el \texttt{code\_challenge}.
  \item Envía solo el \texttt{code\_challenge} al servidor de autorización al
        iniciar el flujo.
  \item Al canjear el código de autorización por el token, el cliente revela el
        \texttt{code\_verifier} original.
  \item El servidor verifica que \texttt{SHA-256(code\_verifier) == code\_challenge}.
\end{enumerate}

Un atacante que intercepte el código de autorización no conoce el
\texttt{code\_verifier} almacenado en la memoria del proceso legítimo, por lo que
el intento de canje fracasará inevitablemente. PKCE es el estándar obligatorio en
2025 para todo tipo de clientes, incluyendo aplicaciones de escritorio
\cite{hypr2025,srivathsav2025}.

\subsection{Niveles de Garantía de Autenticación: NIST SP 800-63-4}
\label{sec:nist}

El NIST publicó en 2025 la cuarta revisión de sus Directrices de Identidad
Digital (SP 800-63-4) \cite{nist80063}, estableciendo tres niveles de garantía
(AAL, \textit{Authenticator Assurance Levels}):

\begin{table}[h!]
\centering
\caption{Niveles de Garantía de Autenticación según NIST SP 800-63-4}
\label{tab:aal}
\begin{tabularx}{\textwidth}{>{\bfseries}l X X}
\toprule
Nivel & Requisitos & Aplicaciones Típicas \\
\midrule
AAL1 & Un solo factor. Desaconsejado para datos personales o sensibles en 2025. &
      Aplicaciones públicas de bajo riesgo. \\[4pt]
AAL2 & MFA obligatoria. \textbf{Passkeys sincronizables autorizados
       explícitamente.} Resistencia al phishing promovida. &
      Nueva línea base para aplicaciones corporativas y comerciales. \\[4pt]
AAL3 & Máxima garantía. \textbf{Prohíbe passkeys sincronizables.} Exige claves
       ligadas al hardware (\textit{device-bound keys}) no exportables (TPM,
       Secure Enclave). &
      Infraestructura gubernamental, sistemas financieros críticos. \\
\bottomrule
\end{tabularx}
\end{table}

La distinción clave del AAL3 es que ni siquiera el proveedor del sistema operativo
puede exportar la clave privada del usuario; ésta queda sellada en el chip
criptográfico del dispositivo físico \cite{nist80063,hypr2025}.

\subsection{Observabilidad Profunda vs. Monitoreo Tradicional}
\label{sec:obs}

\subsubsection{El Abismo Semántico}

El monitoreo tradicional opera sobre métricas de infraestructura predefinidas
(CPU, memoria, latencia) y emite alertas cuando se supera un umbral estático.
Es inherentemente \textbf{reactivo}: detecta ``qué componente falló'' solo ante
problemas conocidos (\textit{known unknowns}). Cuando un atacante usa
credenciales válidas pero robadas, los indicadores de salud permanecen normales.
OWASP cataloga esta incapacidad como A09:2021 \cite{owasptop10}.

La \textbf{observabilidad} es una propiedad arquitectónica diseñada desde la
concepción del sistema: permite inferir su estado interno a partir de la
correlación de sus salidas externas. Es proactiva y causal --- responde
``\textit{por qué falló}'' y permite descubrir anomalías desconocidas
(\textit{unknown unknowns}) \cite{newrelic2025,crowdstrike2025,dynatrace2025}.

\begin{table}[h!]
\centering
\caption{Monitoreo Tradicional vs. Observabilidad Profunda}
\label{tab:monvsobs}
\begin{tabularx}{\textwidth}{>{\bfseries}p{3cm} X X}
\toprule
Dimensión & Monitoreo Tradicional & Observabilidad Profunda (OTel + eBPF) \\
\midrule
Enfoque operacional & Reactivo. ``Qué ocurrió'' / ``El sistema falló''. &
  Proactivo y causal. ``Por qué falló''. Análisis de causa raíz. \\[4pt]
Alcance & Métricas de infraestructura predefinidas y logs locales aislados. &
  Trazas distribuidas correlacionadas + captura a nivel de kernel. \\[4pt]
Resolución de problemas & Fallos conocidos y recurrentes (\textit{knowns}). &
  Fallos complejos y vectores de ataque novedosos (\textit{unknown unknowns}). \\[4pt]
Aplicación en ciberseguridad & Detección de anomalías volumétricas obvias
  (DDoS, bloqueos IP). &
  Análisis comportamental en microsegundos, detección de exfiltración,
  perfilado de atacantes persistentes. \\
\bottomrule
\end{tabularx}
\end{table}

\subsubsection{Los Cuatro Pilares: Marco MELT}

La observabilidad en seguridad de aplicaciones se fundamenta en el marco
\textbf{MELT} (\textit{Metrics, Events, Logs, Traces}) \cite{newrelic2025}:
métricas cuantitativas del sistema, eventos discretos con contexto de negocio,
registros cronológicos estructurados y trazas distribuidas que correlacionan
la ejecución de una solicitud a través de múltiples servicios.

\subsubsection{OpenTelemetry (OTel)}

OpenTelemetry \cite{opentelemetry} es un marco de código abierto respaldado por
la CNCF (\textit{Cloud Native Computing Foundation}) que provee APIs, SDKs y
agentes para la recolección estandarizada y agnóstica al proveedor
(\textit{vendor-neutral}) de datos telemétricos. Su innovación crítica en
ciberseguridad reside en la \textbf{trazabilidad distribuida}:

Al registrarse una solicitud de inicio de sesión, el API Gateway inyecta
automáticamente un \textbf{TraceID} único. Este identificador se propaga en las
cabeceras de red a través de cada microservicio, base de datos y componente de
infraestructura participante, permitiendo reconstruir la historia cronológica
y causal de toda la transacción \cite{opentelemetry,horovits2025}. Si ese mismo
TraceID aparece posteriormente asociado a consultas masivas y atípicas a bases
de datos sensibles, la plataforma puede elevar automáticamente una alerta de
exfiltración.

\subsubsection{eBPF (\textit{Extended Berkeley Packet Filter})}

eBPF \cite{linuxfoundation2025} es una tecnología revolucionaria integrada de
forma nativa en el kernel Linux (y en adopción acelerada en Windows). A diferencia
de los agentes tradicionales que operan en el \textit{user space} --- induciendo
costosos cambios de contexto (\textit{context switches}) ---, eBPF permite
inyectar y ejecutar programas personalizados directamente en el núcleo del
sistema operativo sin recompilar el kernel ni cargar módulos inestables. Antes
de su ejecución, estos programas son verificados por una máquina virtual interna
que garantiza la seguridad del kernel.

Desde la perspectiva de seguridad, eBPF ofrece tres capacidades estratégicas
\cite{researchgate2024,groundcover2025}:

\begin{enumerate}[noitemsep]
  \item \textbf{Instrumentación de Cero Código} (\textit{Zero-Code Instrumentation}):
        Captura llamadas al sistema (syscalls), eventos de red y operaciones de
        disco con sobrecarga prácticamente nula (hasta 210 veces menos que
        agentes en user space), sin requerir que los desarrolladores modifiquen
        el código fuente.
  \item \textbf{Detección de intrusiones a nivel de kernel:} Intercepta
        en microsegundos intentos de lanzar shells remotas, modificar archivos
        críticos del sistema o inyectar backdoors, exportando telemetría al SIEM
        antes de que la instrucción maliciosa se consuma.
  \item \textbf{Inspección pre/post-cifrado TLS:} Al operar antes de que los
        paquetes sean cifrados por TLS o después de que sean descifrados,
        eBPF garantiza que ninguna carga útil maliciosa pueda ocultarse dentro
        de canales aparentemente legítimos y cifrados.
\end{enumerate}

% ═════════════════════════════════════════════════════════════════════════════
\section{Marco Legal Colombiano}
% ═════════════════════════════════════════════════════════════════════════════

El diseño de sistemas de autenticación y observabilidad no es exclusivamente
un problema de ingeniería; constituye una responsabilidad jurídica.

\subsection{Ley 1273 de 2009 --- Delitos Informáticos}

Esta ley modificó el Código Penal colombiano para tipificar los delitos
informáticos, creando el bien jurídico tutelado ``de la protección de la
información y de los datos'' \cite{ley1273}. Los artículos más relevantes para
el diseño de sistemas de autenticación son:

\begin{itemize}[noitemsep]
  \item \textbf{Art. 269A (Acceso abusivo):} Sanciona el acceso a sistemas
        informáticos sin autorización o excediendo los límites acordados. Esto
        exige que la autenticación esté ligada a un sistema robusto de
        autorización (e.g., RBAC) que impida la escalada de privilegios.
  \item \textbf{Art. 269G (Suplantación de sitios web):} Penaliza la creación
        de páginas fraudulentas para capturar credenciales. La implementación
        de passkeys responde directamente a esta amenaza al vincular
        criptográficamente la autenticación al dominio real del servicio.
  \item \textbf{Art. 269I (Hurto por medios informáticos):} Sanciona suplantar
        a un usuario ante sistemas de autenticación para cometer fraude.
\end{itemize}

\subsection{Ley 1581 de 2012 --- Protección de Datos Personales}

Esta ley estatutaria regula el derecho fundamental al \textit{Habeas Data}
\cite{ley1581}. Sus principios rectores imponen obligaciones directas sobre los
sistemas de autenticación y observabilidad:

\begin{itemize}[noitemsep]
  \item \textbf{Principio de finalidad:} Los datos telemétricos recopilados
        (IPs, identificadores de dispositivo, patrones de comportamiento) deben
        usarse exclusivamente para los fines de seguridad informados al usuario.
  \item \textbf{Principio de seguridad:} Obliga a implementar medidas técnicas
        para evitar adulteración, pérdida o acceso no autorizado. En el contexto
        de inicio de sesión, aplica a huellas biométricas, tokens de sesión y
        telemetría comportamental.
  \item \textbf{Principio de libertad:} El tratamiento de datos requiere
        consentimiento previo, expreso e informado del titular.
\end{itemize}

% ═════════════════════════════════════════════════════════════════════════════
\section{Propuesta de Arquitectura Integrada}
% ═════════════════════════════════════════════════════════════════════════════

La propuesta adopta un modelo de \textbf{Seguridad por Diseño} que adapta el
flujo de autenticación a las particularidades de cada plataforma, integrando
una malla de observabilidad profunda en el backend.

\subsection{Implementación por Plataforma}

\subsubsection{Aplicaciones Web}

Las aplicaciones web operan en un entorno inherentemente hostil donde se ejecuta
código JavaScript potencialmente no confiable. Las vulnerabilidades predominantes
son XSS (\textit{Cross-Site Scripting}) y CSRF
(\textit{Cross-Site Request Forgery}) \cite{stackhawk2025}. Las estrategias
recomendadas son:

\begin{itemize}[noitemsep]
  \item \textbf{WebAuthn para passkeys:} Registrar los dispositivos del usuario
        (sensores biométricos integrados, llaves FIDO USB) como autenticadores
        de primera línea.
  \item \textbf{Patrón BFF y cookies HttpOnly:} En lugar de almacenar JWT en
        \texttt{localStorage} (vulnerable a XSS), utilizar el patrón
        \textit{Backend-for-Frontend} o cookies con atributos
        \texttt{HttpOnly}, \texttt{Secure} y \texttt{SameSite=Strict}
        \cite{srivathsav2025}.
  \item \textbf{Cabeceras de seguridad HTTP:} TLS 1.3 obligatorio, HSTS
        (\textit{HTTP Strict Transport Security}) y una política CSP
        (\textit{Content Security Policy}) restrictiva que controle el origen
        de scripts ejecutables \cite{oligo2025}.
\end{itemize}

\subsubsection{Aplicaciones Móviles Nativas (iOS y Android)}

Los dispositivos móviles combinan conectividad persistente con exposición física
(redes Wi-Fi públicas, posibilidad de robo). El OWASP Mobile Top 10
\cite{owaspMobile} subraya el almacenamiento inseguro de datos como
vulnerabilidad crítica. Las contramedidas son:

\begin{itemize}[noitemsep]
  \item \textbf{Almacenamiento respaldado por hardware:} Los \textit{refresh
        tokens} deben residir exclusivamente en el \textbf{Keychain} (iOS/Secure
        Enclave) o el \textbf{Keystore} (Android/TEE o Titan M), garantizando
        que no puedan extraerse en texto plano incluso ante un rooting o
        jailbreak completo \cite{mdpi2025}.
  \item \textbf{Autenticación biométrica como precondición:} FaceID, TouchID o
        reconocimiento de iris deben requerirse antes de desbloquear
        criptográficamente los tokens almacenados.
  \item \textbf{Certificate Pinning:} La aplicación rechaza conexiones TLS a
        menos que el certificado del servidor coincida con un hash embebido en
        el binario compilado, neutralizando ataques AiTM en redes comprometidas
        \cite{mdpi2025,touchlane2025}.
\end{itemize}

\subsubsection{Aplicaciones de Escritorio Nativas}

Las aplicaciones de escritorio son especialmente vulnerables a la ingeniería
inversa, la inyección en DLLs y el volcado de memoria (herramientas como
Mimikatz) \cite{chudovo2025}. Las estrategias son:

\begin{itemize}[noitemsep]
  \item \textbf{DPAPI / Keychain Access:} En Windows, la Data Protection API
        (DPAPI) cifra los tokens vinculándolos a la entropía de las credenciales
        de sesión local del usuario. En macOS, se usa Keychain Access.
  \item \textbf{Credential Guard:} En despliegues corporativos (Windows
        Server 2025 / Windows 11 Enterprise), esta característica de
        virtualización aísla los secretos de autenticación en un contenedor
        de hipervisor, inaccesible incluso para procesos con privilegios
        SYSTEM \cite{microsoft2025}.
  \item \textbf{OAuth 2.0 + PKCE para Custom URI Schemes:} Las aplicaciones
        de escritorio que utilizan manejadores de URL personalizados para recibir
        el código de redirección deben exigir PKCE para impedir que malware
        registrado en el mismo sistema secuestre el código de autorización
        \cite{srivathsav2025}.
\end{itemize}

\subsection{Pipeline de Observabilidad para Detección Activa de Amenazas}

\subsubsection{OpenTelemetry: Trazabilidad Estructural del Negocio}

Se instrumentan los microservicios IAM (\textit{Identity and Access Management})
con los SDKs de OTel. El flujo propuesto es:

\begin{enumerate}[noitemsep]
  \item El API Gateway genera un \textbf{TraceID} global al recibir cualquier
        solicitud (web, móvil o escritorio).
  \item El TraceID se propaga en cabeceras HTTP/gRPC por cada componente
        participante (microservicios, bases de datos, colas de mensajes).
  \item El equipo SecOps puede rastrear con precisión forense la ruta temporal
        y causal de cualquier acción post-autenticación.
  \item Si el TraceID de una sesión legítima aparece en consultas masivas a
        datos sensibles, la plataforma eleva automáticamente una alerta de
        posible exfiltración \cite{opentelemetry,horovits2025,thenewstack2025}.
\end{enumerate}

\subsubsection{eBPF: Sensores de Seguridad en Tiempo de Ejecución}

Los sensores eBPF (implementaciones como OBI --- \textit{OpenTelemetry eBPF
Instrumentation} \cite{obiRelease} o Tetragon) se adjuntan a
\textit{tracepoints} y \textit{kprobes} específicos del kernel, observando:

\begin{itemize}[noitemsep]
  \item Cada conexión TCP iniciada por los contenedores de autenticación.
  \item Cada resolución DNS y llamada al sistema (syscall) de los procesos
        involucrados.
  \item Intentos de lanzar shells interactivas, modificar archivos críticos del
        sistema o instalar backdoors.
\end{itemize}

Esta aproximación garantiza cobertura telemétrica incluso en componentes de
código cerrado o infraestructuras \textit{legacy}, sobre las que el equipo
de ingeniería no tiene acceso al código fuente \cite{groundcover2025,
linuxfoundation2025}.

\subsubsection{Prevención de ATO mediante Analítica e Inteligencia Artificial}

La telemetría unificada de OTel y eBPF alimenta motores de analítica
comportamental e IA, habilitando dos capacidades avanzadas \cite{clear2025,
digitalocean2025}:

\begin{enumerate}[noitemsep]
  \item \textbf{Detección de Credential Stuffing en Tiempo Real:} El motor
        detecta que cientos de intentos de inicio de sesión, aunque provienen
        de IPs globalmente diversas, comparten patrones sutiles en cabeceras
        HTTP, temporización y rastros de red, revelándolos como botnets
        coordinadas. El WAF levanta barreras dinámicas automáticamente.
  \item \textbf{Biometría Comportamental Continua:} Durante la sesión activa
        (no solo en el inicio de sesión), la IA analiza la dinámica del cursor,
        la cadencia de escritura y el ángulo de sostenimiento del dispositivo
        para perfilar al usuario legítimo. Si el comportamiento detectado
        corresponde a un script automatizado o a un intruso remoto, el sistema
        revoca inmediatamente los JWT asociados y exige re-autenticación
        biométrica.
\end{enumerate}

% ═════════════════════════════════════════════════════════════════════════════
\section{Evaluación y Conformidad de la Arquitectura}
% ═════════════════════════════════════════════════════════════════════════════

\subsection{Mitigación del OWASP Top 10}

La arquitectura propuesta aborda directamente las categorías OWASP más críticas
\cite{owasptop10,owaspSBD}:

\begin{itemize}[noitemsep]
  \item \textbf{A01:2021 --- Broken Access Control:} Mitigado mediante RBAC
        integrado con el flujo OAuth 2.0/OIDC y políticas JWT granulares.
  \item \textbf{A07:2021 --- Identification and Authentication Failures:}
        Erradicado algorítmicamente con passkeys FIDO2/WebAuthn, que eliminan
        las contraseñas almacenadas en servidores.
  \item \textbf{A09:2021 --- Security Logging and Monitoring Failures:}
        Resuelto mediante la malla OTel + eBPF, que garantiza cobertura
        telemétrica sin puntos ciegos, incluso en arquitecturas de contenedores
        efímeros.
\end{itemize}

\subsection{Alineación con NIST SP 800-63-4}

La arquitectura es plenamente compatible con la revisión 2025 del NIST
\cite{nist80063}:

\begin{itemize}[noitemsep]
  \item \textbf{AAL2:} Passkeys sincronizables (iCloud Keychain, Google
        Password Manager) satisfacen los requisitos de MFA resistente al
        phishing para aplicaciones comerciales estándar.
  \item \textbf{AAL3:} Claves ligadas al hardware (TPM, Secure Enclave)
        satisfacen el requisito de \textit{device-bound keys} no exportables
        para infraestructura crítica.
\end{itemize}

\subsection{Cumplimiento del Marco Legislativo Colombiano}

\begin{itemize}[noitemsep]
  \item \textbf{Ley 1273/2009 \cite{ley1273}:} La telemetría capturada por
        eBPF genera evidencia digital forense (cadena de custodia) admisible
        ante la Fiscalía General de la Nación y el CAI Virtual de la Policía,
        facilitando la persecución del acceso abusivo (Art. 269A).
  \item \textbf{Ley 1581/2012 \cite{ley1581}:} La arquitectura FIDO2 elimina
        el almacenamiento de contraseñas y secretos biométricos en servidores
        centralizados, reduciendo drásticamente el riesgo de violación del
        principio de seguridad del Habeas Data. La telemetría de observabilidad
        se limita a fines de seguridad, en cumplimiento del principio de finalidad.
\end{itemize}

% ═════════════════════════════════════════════════════════════════════════════
\section{Conclusiones}
% ═════════════════════════════════════════════════════════════════════════════

La ciberseguridad de aplicaciones en 2025 no puede fundamentarse en la premisa
de que el perímetro de red protege los sistemas internos. Con la disolución de
las fronteras tradicionales, el proceso de inicio de sesión y los controles de
autorización se han convertido en el epicentro del control de acceso, principio
rector de cualquier implementación de Arquitectura Zero Trust.

El análisis exhaustivo presentado demuestra que la eliminación de la dependencia
de contraseñas estáticas es un mandato técnico y regulatorio. Los passkeys,
fundamentados en FIDO2/WebAuthn, erradican algorítmicamente el credential stuffing
y el phishing al nivel del protocolo, ofrecen tasas de éxito superiores al 98\%
y son compatibles con los niveles AAL2 y AAL3 del NIST SP 800-63-4. La extensión
PKCE en flujos OAuth 2.0 es igualmente obligatoria para todos los clientes públicos,
garantizando que la interceptación del código de autorización sea operativamente
inútil para un atacante.

Sin embargo, un flujo de autenticación robusto es solo el primer eslabón de la
cadena. La integración estructural de OpenTelemetry y eBPF transforma la postura
de seguridad de reactiva a proactiva: OTel teje una trazabilidad distribuida que
correlaciona la identidad del usuario con cada acción post-autenticación, mientras
que eBPF provee visibilidad a nivel de kernel sin alterar el código fuente ni
introducir latencia apreciable. Juntos, habilitan la detección de exfiltración,
movimiento lateral y ataques persistentes que son completamente invisibles para
el monitoreo tradicional.

Finalmente, la arquitectura satisface simultáneamente los requisitos del OWASP
Top 10, el NIST SP 800-63-4 y el marco jurídico colombiano (Leyes 1273 de 2009
y 1581 de 2012), demostrando que la seguridad técnica y el cumplimiento normativo
no son objetivos contrapuestos sino sinérgicos. La conjunción de gestión de
identidades modernas y observabilidad omnipresente no debe entenderse como una
tendencia tecnológica aislada, sino como la única estrategia arquitectónica
viable para la supervivencia digital en la era hiperconectada.

% ═════════════════════════════════════════════════════════════════════════════
\newpage
\bibliographystyle{IEEEtran}
\begin{thebibliography}{99}

\bibitem{authx2025}
AuthX, ``Account Takeover Statistics 2025: Key Trends \& Insights,''
\textit{AuthX Blog}, 2025. [Online]. Available:
\url{https://www.authx.com/blog/account-takeover-fraud-statistics-insights/}

\bibitem{verizon2025}
Verizon, ``Additional 2025 DBIR Research on Credential Stuffing,''
\textit{Verizon Business Resources}, 2025. [Online]. Available:
\url{https://www.verizon.com/business/resources/articles/credential-stuffing-attacks-2025-dbir-research/}

\bibitem{mojoauth2025}
MojoAuth, ``Passkeys Handbook 2025: Secure, Passwordless Authentication,''
\textit{MojoAuth White Papers}, 2025. [Online]. Available:
\url{https://mojoauth.com/white-papers/passkeys-passwordless-authentication-handbook/}

\bibitem{nist80063}
NIST, ``Digital Identity Guidelines: SP 800-63-4,'' \textit{NIST Special
Publication 800-63-4}, National Institute of Standards and Technology, 2025.
[Online]. Available: \url{https://pages.nist.gov/800-63-4/sp800-63.html}

\bibitem{nistZTA}
NIST NCCoE, ``Implementing a Zero Trust Architecture,'' \textit{NIST
National Cybersecurity Center of Excellence}, 2025. [Online]. Available:
\url{https://www.nccoe.nist.gov/projects/implementing-zero-trust-architecture}

\bibitem{owasptop10}
OWASP Foundation, ``OWASP Top Ten 2021,'' \textit{Open Web Application
Security Project}, 2021. [Online]. Available:
\url{https://owasp.org/www-project-top-ten/}

\bibitem{owaspMobile}
OWASP Foundation, ``Mobile Top 10 2024,'' \textit{Open Web Application
Security Project}, 2024. [Online]. Available:
\url{https://owasp.org/www-project-mobile-top-10/}

\bibitem{owaspSBD}
OWASP Foundation, ``OWASP Secure by Design Framework,'' \textit{Open Web
Application Security Project}, 2024. [Online]. Available:
\url{https://owasp.org/www-project-secure-by-design-framework/}

\bibitem{opentelemetry}
OpenTelemetry, ``OpenTelemetry: An Observability Framework for Cloud-Native
Software,'' \textit{Cloud Native Computing Foundation}, 2025. [Online].
Available: \url{https://opentelemetry.io/}

\bibitem{horovits2025}
D. Horovits, ``Unveiling OpenTelemetry eBPF Instrumentation,''
\textit{Medium}, 2025. [Online]. Available:
\url{https://horovits.medium.com/unveiling-opentelemetry-ebpf-instrumentation-377cb0432bf1}

\bibitem{obiRelease}
OpenTelemetry, ``OpenTelemetry eBPF Instrumentation Marks the First
Release,'' \textit{OpenTelemetry Blog}, 2025. [Online]. Available:
\url{https://opentelemetry.io/blog/2025/obi-announcing-first-release/}

\bibitem{linuxfoundation2025}
Linux Foundation, ``eBPF for the Infrastructure Platform: The State of
eBPF 2025,'' \textit{Linux Foundation}, 2025. [Online]. Available:
\url{https://www.linuxfoundation.org/hubfs/eBPF/The_State_of_eBPF25_111925.pdf}

\bibitem{researchgate2024}
J. Doe \textit{et al.}, ``eBPF-Enhanced Complete Observability Solution for
Cloud-Native Microservices,'' \textit{ResearchGate}, 2024. [Online].
Available:
\url{https://www.researchgate.net/publication/379891642}

\bibitem{groundcover2025}
Groundcover, ``OpenTelemetry and eBPF: Everything You Need to Know,''
\textit{Groundcover Blog}, 2025. [Online]. Available:
\url{https://www.groundcover.com/ebpf/opentelemetry-ebpf}

\bibitem{frontegg2025}
Frontegg, ``What Is User Authentication? A 2025 Guide for Modern Apps,''
\textit{Frontegg Blog}, 2025. [Online]. Available:
\url{https://frontegg.com/blog/authentication}

\bibitem{srivathsav2025}
J. R. S. Adari, ``Building Secure Authentication Layers: A Complete Guide
for 2025,'' \textit{srivathsav.me}, 2025. [Online]. Available:
\url{https://srivathsav.me/blog/building-secure-authentication-layers-2025}

\bibitem{hypr2025}
HYPR, ``NIST SP 800-63-3 \& 63-4: Digital Identity Guidelines Review,''
\textit{HYPR Blog}, 2025. [Online]. Available:
\url{https://www.hypr.com/blog/nist-sp-800-63-3-digital-identity-guidelines-review}

\bibitem{authsignal2025}
Authsignal, ``Passwordless Authentication in 2025: The Year Passkeys Went
Mainstream,'' \textit{Authsignal Blog}, 2025. [Online]. Available:
\url{https://www.authsignal.com/blog/articles/passwordless-authentication-in-2025-the-year-passkeys-went-mainstream}

\bibitem{googleModernPass}
Google Cloud, ``Modern Password Security for System Designers,''
\textit{Google Cloud Documentation}, 2024. [Online]. Available:
\url{https://docs.cloud.google.com/solutions/modern-password-security-for-system-designers.pdf}

\bibitem{newrelic2025}
New Relic, ``Observability vs. Monitoring: What's the Difference?''
\textit{New Relic Blog}, 2025. [Online]. Available:
\url{https://newrelic.com/blog/observability/observability-vs-monitoring}

\bibitem{crowdstrike2025}
CrowdStrike, ``Observability vs. Monitoring: What's the Difference?''
\textit{CrowdStrike Cybersecurity 101}, 2025. [Online]. Available:
\url{https://www.crowdstrike.com/en-us/cybersecurity-101/observability/observability-vs-monitoring/}

\bibitem{dynatrace2025}
Dynatrace, ``Observability vs. Monitoring: What's the Difference?''
\textit{Dynatrace Blog}, 2025. [Online]. Available:
\url{https://www.dynatrace.com/news/blog/observability-vs-monitoring/}

\bibitem{thenewstack2025}
The New Stack, ``Observability in 2025: OpenTelemetry and AI to Fill In
Gaps,'' \textit{The New Stack}, 2025. [Online]. Available:
\url{https://thenewstack.io/observability-in-2025-opentelemetry-and-ai-to-fill-in-gaps/}

\bibitem{cncf2025}
CNCF, ``Observability Trends in 2025 -- What's Driving Change?''
\textit{Cloud Native Computing Foundation Blog}, 2025. [Online]. Available:
\url{https://www.cncf.io/blog/2025/03/05/observability-trends-in-2025-whats-driving-change/}

\bibitem{mdpi2025}
T. Al-Turjman \textit{et al.}, ``A Review on Secure Authentication
Mechanisms for Mobile Security,'' \textit{MDPI Sensors}, vol. 25, no. 3,
p. 700, 2025. DOI: 10.3390/s25030700.

\bibitem{touchlane2025}
Touchlane, ``Common Mobile Application Security Vulnerabilities 2025,''
\textit{Touchlane Blog}, 2025. [Online]. Available:
\url{https://touchlane.com/common-mobile-application-vulnerabilities-2025/}

\bibitem{microsoft2025}
Microsoft, ``What's New in Windows Server 2025,'' \textit{Microsoft Learn},
2025. [Online]. Available:
\url{https://learn.microsoft.com/en-us/windows-server/get-started/whats-new-windows-server-2025}

\bibitem{chudovo2025}
Chudovo, ``Why Desktop Applications Will Still Be in Demand in 2025,''
\textit{Chudovo Blog}, 2025. [Online]. Available:
\url{https://chudovo.com/why-desktop-applications-will-still-be-in-demand/}

\bibitem{stackhawk2025}
StackHawk, ``10 Web Application Security Threats for 2025 \& How to
Respond,'' \textit{StackHawk Blog}, 2025. [Online]. Available:
\url{https://www.stackhawk.com/blog/10-web-application-security-threats-and-how-to-mitigate-them/}

\bibitem{oligo2025}
Oligo Security, ``Web App Security in 2025: 5 Technologies You Can't Do
Without,'' \textit{Oligo Academy}, 2025. [Online]. Available:
\url{https://www.oligo.security/academy/web-app-security-in-2025-5-technologies-you-cant-do-without}

\bibitem{zerotrust2025}
Northern Technologies Group, ``Zero Trust Architecture in 2025: Shifting
from Perimeter Security to `Never Trust, Always Verify',''
\textit{NTG Blog}, 2025. [Online]. Available:
\url{https://ntgit.com/zero-trust-architecture-in-2025-shifting-from-perimeter-security-to-never-trust-always-verify/}

\bibitem{clear2025}
CLEAR, ``Fraud Analytics in 2025: Why Traditional Systems Fail,''
\textit{CLEAR Identity Blog}, 2025. [Online]. Available:
\url{https://identity.clearme.com/post/fraud-analytics-2025-why-traditional-systems-fail}

\bibitem{digitalocean2025}
DigitalOcean, ``Understanding AI Fraud Detection and Prevention in 2026,''
\textit{DigitalOcean Resources}, 2025. [Online]. Available:
\url{https://www.digitalocean.com/resources/articles/ai-fraud-detection}

\bibitem{acis2024}
ACIS, ``Colombia enfrenta récord de 36.000 millones de ciberataques en
2024,'' \textit{Asociación Colombiana de Ingenieros de Sistemas}, 2024.
[Online]. Available:
\url{https://www.acis.org.co/blog/noticias-2/colombia-enfrenta-record-de-36000-millones-de-ciberataques-en-2024-4964}

\bibitem{ccit2024}
CCIT, ``Balance de Ciberseguridad 2024: Desafíos y prevención para un
entorno digital seguro,'' \textit{Cámara Colombiana de Informática y
Telecomunicaciones}, 2024. [Online]. Available:
\url{https://www.ccit.org.co/noticias/balance-de-ciberseguridad-2024-desafios-y-prevencion-para-un-entorno-digital-seguro/}

\bibitem{ley1273}
Congreso de la República de Colombia, ``Ley 1273 de 2009: Por medio de la
cual se modifica el Código Penal, se crea un nuevo bien jurídico tutelado
denominado `de la protección de la información y de los datos','' 2009.
[Online]. Available:
\url{https://www.armada.mil.co/sites/default/files/normograma_arc/telematica/Ley\%201273\%20de\%202009\%20Modifica\%20CP\%20para\%20protecci\%C3\%B3n\%20de\%20datos\%20e\%20informaci\%C3\%B3n.pdf}

\bibitem{ley1581}
Congreso de la República de Colombia, ``Ley 1581 de 2012: Por la cual se
dictan disposiciones generales para la protección de datos personales,''
2012. [Online]. Available:
\url{https://www.alcaldiabogota.gov.co/sisjur/normas/Norma1.jsp?i=49981}

\bibitem{iso27001}
International Organization for Standardization, ``ISO/IEC 27001:2022 --
Information Security Management Systems,'' \textit{ISO}, 2022.

\bibitem{authsignal_stuffing}
Authsignal, ``How to Actually Stop Credential Stuffing in 2025,''
\textit{Authsignal Blog}, 2025. [Online]. Available:
\url{https://www.authsignal.com/blog/articles/how-to-actually-stop-credential-stuffing-in-2025}

\end{thebibliography}

\end{document}
